\documentclass[11pt]{article}
\usepackage[utf8]{inputenc}
\usepackage[T1]{fontenc}
\usepackage{times}
\usepackage{amsmath,amssymb}
\usepackage{booktabs}
\usepackage{multirow}
\usepackage{graphicx}
\usepackage{xcolor}
\usepackage{hyperref}
\usepackage{url}
\usepackage[margin=1in]{geometry}
\usepackage{natbib}
\usepackage{subcaption}
\usepackage{tablefootnote}

\newcommand{\gar}{\textsc{GAR}}
\newcommand{\garours}{\textsc{GAR}\textsubscript{ours}}
\newcommand{\gardongre}{\textsc{GAR}\textsubscript{Dongre}}
\newcommand{\tg}{\textsc{ThreadGate}}
\newcommand{\ug}{\textsc{UniformGate}}
\newcommand{\lo}{\textsc{LoRA-only}}
\renewcommand{\ss}{\textsc{SS}}

\title{From Collapse to Recovery: Diagnosing and Mitigating\\ Exposure Bias in Multi-Turn Attention Objectives}

\author{Anonymous}
\date{}

\begin{document}
\maketitle

% ============================================================
% ABSTRACT
% ============================================================
\begin{abstract}
Multi-turn dialogue models suffer progressive goal-attention degradation as conversations grow longer. Prior work shows goal information persists in residual streams (probe AUC$=$0.98) even as attention access declines, motivating learnable residual gating to recover it. We implement this hypothesis as \tg{}, a per-layer gating module trained jointly with next-token prediction (NTP) and an auxiliary Goal Accessibility Ratio (\gar{}) loss. \tg{} successfully optimizes \gar{} under teacher-forcing (TF \gar{}$=$0.50), yet collapses at inference on model-generated text (AR \gar{}$=$0.064 on Qwen-2.5-7B, 0.030 on Llama-3.1-8B), revealing an 8$\times$--27$\times$ training--inference gap. Through PCGrad causal ablation, per-head KL divergence analysis, and cross-model replication, we identify exposure bias---not gradient competition---as the root cause. This gap selectively disrupts summary-position attention patterns while leaving all-token attention unaffected. Scheduled sampling (SS) compresses the gap from 8$\times$ to ${\sim}$1.3$\times$ on both Qwen and Llama (3-seed AR \gar{}$=$0.40$\pm$0.05, no perplexity cost), and inference-time $\beta$-steering further recovers \gar{} to 0.72. Our results provide the first empirical characterization of how exposure bias destroys auxiliary attention objectives and demonstrate that classical mitigation transfers effectively to multi-turn transformer settings.
\end{abstract}

% ============================================================
% 1. INTRODUCTION
% ============================================================
\section{Introduction}

Large Language Models (LLMs) exhibit progressive degradation in multi-turn dialogue: as conversations grow longer, models increasingly fail to maintain attention to the user's original goals and instructions \citep{dongre2026goal,liu2024lost,hsieh2024ruler}. \citet{dongre2026goal} formalized this through the Goal Accessibility Ratio (\gar{}), demonstrating 27--48\% declines across 12 models over 50 conversational turns. This degradation directly impacts practical deployments where sustained goal-tracking is essential---customer service, tutoring, and collaborative problem-solving.

A striking observation from \citet{dongre2026goal} is that goal information \emph{persists} in the residual stream even as the attention channel closes: linear probes achieve AUC$=$0.98 on goal-presence classification at late layers. This creates a natural hypothesis: if goal-relevant information survives in residual representations but attention fails to route it, a learnable gate on the residual connection should modulate layer contributions to recover goal accessibility.

We implement this hypothesis as \tg{}, a per-layer gating module that receives hidden states and turn-position embeddings, producing per-token, per-layer gate values trained jointly with NTP and a \gar{} auxiliary loss. Under teacher-forcing, \tg{} raises \gar{} from 0.02 to 0.50---a successful optimization. However, when evaluated on model-generated text (the autoregressive regime), \gar{} collapses to 0.064 on Qwen-2.5-7B (7.8$\times$ gap) and 0.030 on Llama-3.1-8B (15.9$\times$ gap). This dramatic training--inference discrepancy is our central finding.

Through systematic causal analysis, we establish that exposure bias---the distribution shift between gold training contexts and model-generated inference contexts---is the sole root cause of this collapse. PCGrad \citep{yu2020gradient} ablation eliminates gradient direction conflicts between NTP and \gar{} losses without reducing the gap (modus tollens). A $\lambda_{\text{GAR}}$ sweep across two orders of magnitude (0.1--10) eliminates gradient magnitude imbalance without reducing the gap. Per-head KL divergence analysis reveals that the heads most disrupted by distribution shift ($r=0.38$, $p=0.002$) are precisely those responsible for goal-relevant attention routing. Cross-model replication on Llama-3.1-8B confirms the pattern is architecture-general.

Crucially, we demonstrate that scheduled sampling (SS)---originally designed for RNN sequence models \citep{bengio2015scheduled}---transfers effectively to multi-turn transformer attention objectives. SS compresses the TF--AR gap from 8$\times$ to ${\sim}$1.3$\times$ on Qwen-2.5-7B (3-seed: AR \gar{}$=$0.401$\pm$0.047) and from 15.9$\times$ to 1.34$\times$ on Llama-3.1-8B, with no perplexity cost. Inference-time $\beta$-steering independently recovers \gar{} from 0.064 to 0.727. Our contributions are:

\begin{enumerate}
    \item The first empirical characterization of how exposure bias destroys auxiliary attention objectives: an 8$\times$--27$\times$ TF--AR gap specific to summary-position attention patterns (Qwen 7.8$\times$, Llama 15.9$\times$).
    \item Causal exclusion of gradient competition via PCGrad ablation and $\lambda$ sensitivity analysis, isolating exposure bias as the sole mechanism.
    \item Per-head KL divergence analysis localizing the distribution shift to late-layer goal-routing heads ($r=0.38$, $p=0.002$), explaining the selectivity of the collapse.
    \item Cross-model validation on two architectures (Qwen-2.5-7B MHA and Llama-3.1-8B GQA) confirming architecture-general failure and recovery.
    \item Scheduled sampling as an effective mitigation, compressing the gap to ${\sim}$1.3$\times$ on both models (3-seed validated), demonstrating that classical exposure bias remedies transfer to multi-turn attention objectives.
\end{enumerate}

% ============================================================
% 2. RELATED WORK
% ============================================================
\section{Related Work}

\paragraph{Multi-Turn Dialogue Degradation.}
Several studies document performance degradation in long conversations. \citet{dongre2026goal} provide the most systematic treatment through \gar{}, demonstrating progressive goal-attention decay across 12 LLMs. \citet{liu2024lost} identify ``lost-in-the-middle'' phenomena where models struggle with information placed centrally in long contexts. \citet{hsieh2024ruler} introduce RULER for evaluating long-context capabilities, revealing substantial performance drops beyond training lengths. MT-Bench-101 \citep{bai2024mtbench101} and MultiChallenge \citep{li2024multichallenge} provide multi-turn evaluation benchmarks documenting instruction-following decay.

\paragraph{Residual Gating and Adaptive Depth.}
Residual connections are fundamental to deep transformer training \citep{he2016deep}. ReZero \citep{bachlechner2021rezero} initializes residual gates to zero for faster convergence. \citet{raposo2024mixtureofdepths} introduce Mixture-of-Depths for adaptive computation allocation. FlexiDepth \citep{luo2025flexidepth} enables dynamic layer selection. Layer-selective rank reduction \citep{sharma2024truth} shows that not all layers contribute equally. Our per-layer gating follows the ReZero parameterization $x_{l+1} = x_l + g_l \cdot f_l(x_l)$.

\paragraph{Attention Steering and Intervention.}
Attention sinks \citep{xiao2024efficient} reveal fixed attention patterns at sequence boundaries. Gated attention mechanisms \citep{qiu2025gated} learn selective attention routing. Activation steering \citep{turner2023activation,li2024inference} modifies internal representations at inference time to control model behavior. Our $\beta$-steering approach builds on this paradigm.

\paragraph{Exposure Bias.}
The training--inference distribution mismatch in teacher-forced models was formalized by \citet{ranzato2016sequence} and has been extensively studied in machine translation \citep{wiseman2016sequence,zhang2019bridging}. Scheduled sampling \citep{bengio2015scheduled} provides a simple curriculum-based mitigation. Professor Forcing \citep{lamb2016professor} uses adversarial training to align distributions. \citet{dou2022parallel} propose Parallel Attention Forcing for NMT with data-dependent training mode selection. Our work extends exposure bias analysis to \emph{auxiliary} attention losses in multi-turn settings, a previously unexplored combination.

\paragraph{Gradient Conflict in Multi-Task Learning.}
PCGrad \citep{yu2020gradient} projects conflicting gradients to eliminate inter-task interference. GradNorm \citep{chen2018gradnorm} dynamically balances task weights. We use PCGrad not as an optimization strategy but as a \emph{causal ablation} tool: by eliminating gradient conflicts and observing no effect, we rule out gradient competition as a mechanism.

% ============================================================
% 3. BACKGROUND
% ============================================================
\section{Background}

\subsection{Goal Accessibility Ratio}

We adopt the Goal Accessibility Ratio (\gar{}) from \citet{dongre2026goal}, which measures the fraction of a model's attention directed toward goal-relevant tokens. Following Dongre et al., we define two variants:

\textbf{\gardongre{}} uses all response tokens as query positions and system-prompt task tokens as goal positions, enabling direct comparison with \citet{dongre2026goal}.

\textbf{\garours{}} uses the last token of each assistant turn as the query position and system-prompt plus first-user-query tokens as goal positions. The last token serves as a natural summary position whose attention pattern reflects the model's accumulated goal representation.

Both variants are evaluated under two regimes: \textit{teacher-forced} (TF), where the model conditions on gold conversation history, and \textit{autoregressive} (AR), where the model generates its own responses. The TF--AR gap quantifies how well training-time attention patterns generalize to inference.

\subsection{Residual Stream Persistence}

\citet{dongre2026goal} demonstrate that linear probes trained on residual stream activations achieve AUC up to 0.98 for goal-presence classification, even at layers and turn positions where attention-based \gar{} has declined substantially. This dissociation between residual information and attention access provides the empirical motivation for residual gating: the information exists but the attention mechanism fails to access it.

% ============================================================
% 4. METHOD
% ============================================================
\section{Method}

\subsection{ThreadGate Architecture}

\tg{} adds a learnable scalar gate $g_l \in (0,1)$ to each transformer layer $l$, modifying the residual connection as:
\begin{equation}
    x_{l+1} = x_l + g_l \cdot f_l(x_l)
\label{eq:gate}
\end{equation}
where $f_l$ denotes the $l$-th transformer block. When $g_l \to 0$, the layer becomes a no-op (identity); when $g_l \to 1$, the standard residual connection is recovered.

Each gate is computed by a two-layer MLP that takes as input a projection of the hidden state concatenated with turn-position embeddings:
\begin{equation}
    g_l = \sigma\big(W_2 \cdot \text{ReLU}(W_1 \cdot [\text{proj}(x_l);\ \text{emb}(t_d, t_i)] + b_1) + b_2\big)
\end{equation}
where $t_d$ is the distance to the current turn and $t_i$ is the absolute turn index. This design allows gates to be both content-dependent and turn-aware. \ug{} replaces turn embeddings with a constant vector, preserving content-dependence while removing turn awareness. \lo{} fixes $g_l = 1$ for all layers, training only the LoRA adapters with the same \gar{} loss.

\subsection{Training Objective}

The model is trained with a joint loss:
\begin{equation}
    \mathcal{L} = \mathcal{L}_{\text{NTP}} + \lambda_{\text{GAR}} \cdot \mathcal{L}_{\text{GAR}} + \lambda_{\text{div}} \cdot \mathcal{L}_{\text{div}}
\end{equation}

$\mathcal{L}_{\text{GAR}}$ uses an MSE target formulation that trains attention patterns toward a target \gar{} value. $\mathcal{L}_{\text{div}}$ is a gate diversity regularizer that prevents all gates from collapsing to the same value. We arrived at this formulation through systematic exploration of four loss families across seven configurations (v1--v7): cumulative (v1--v2), maximize (v3--v5), smooth saturation (v6), and MSE target (v7). Only the MSE+gate\_div combination achieved stable training convergence.

\subsection{Training Setup}

All models are trained on UltraChat \citep{ding2023enhancing} multi-turn conversations with LoRA adapters (rank 32) on attention projection matrices. Qwen-2.5-7B \citep{yang2024qwen2} serves as the primary model; Llama-3.1-8B \citep{dubey2024llama} provides cross-model validation. Training uses 3 epochs, $\lambda_{\text{GAR}}=1.0$, $\lambda_{\text{div}}=0.01$, and AdamW with learning rate $2 \times 10^{-4}$.

\subsection{Scheduled Sampling}

To mitigate exposure bias, we apply scheduled sampling (SS) during training. At each training step, with probability $p(t)$ that increases linearly from 0 to $p_{\text{max}}$, the model's own generated tokens replace the gold tokens in the conversation context for subsequent turns. This exposes the model to its own distribution during training, reducing the TF--AR distribution gap. For Qwen, $p_{\text{max}} = 0.5$; for Llama, $p_{\text{max}} = 0.3$.

\subsection{Inference-Time $\beta$-Steering}

At inference, we apply a steering coefficient $\beta$ that amplifies gate deviations from identity:
\begin{equation}
    g_l^{\text{steer}} = \sigma\big(\text{logit}(g_l) \cdot (1 + \beta)\big)
\end{equation}
This allows post-hoc control of goal-attention strength without retraining.

% ============================================================
% 5. EXPERIMENTS
% ============================================================
\section{Experiments}

\subsection{Setup}

We evaluate on 200 held-out multi-turn conversations (10 turns each) with bootstrap confidence intervals. AR evaluation generates model responses autoregressively and computes \garours{} on the resulting text. TF evaluation computes \garours{} on gold text. We report \garours{} as the primary metric and perplexity (PPL) on the same conversations as a language modeling quality check.

\subsection{Main Results: Gating Provides Zero Benefit}

Table~\ref{tab:main} presents the four-way ablation on Qwen-2.5-7B. Under teacher-forcing, all fine-tuned variants achieve similar \garours{} (${\sim}$0.50) and PPL (${\sim}$2.78), with bootstrap 95\% CIs for all pairwise differences including zero (Table~\ref{tab:ci}). Under autoregressive generation, \garours{} collapses uniformly across all variants: \tg{} (0.064), \ug{} (0.062), and \lo{} (0.063)---all within noise of each other and barely above the base model (0.019). The 7.8$\times$ TF--AR gap is identical whether gating is learned, uniform, or absent entirely.

\begin{table}[t]
\centering
\small
\caption{Four-way ablation on Qwen-2.5-7B. Gating provides zero incremental benefit over \lo{} in both TF and AR regimes. All fine-tuned variants exhibit the same ${\sim}$8$\times$ TF--AR gap.}
\label{tab:main}
\begin{tabular}{lcccc}
\toprule
\textbf{Method} & \textbf{TF \garours{}} & \textbf{AR \garours{}} & \textbf{Gap} & \textbf{PPL} \\
\midrule
Base & 0.020 & 0.019 & 1.1$\times$ & 2.85 \\
\lo{} & 0.497 & 0.063 & 7.9$\times$ & 2.78 \\
\ug{} & 0.495 & 0.062 & 8.0$\times$ & 2.78 \\
\tg{} (v7) & 0.500 & 0.064 & 7.8$\times$ & 2.78 \\
\tg{} + PCGrad (v9) & 0.497 & 0.061 & 8.1$\times$ & 2.78 \\
\bottomrule
\end{tabular}
\end{table}

\begin{table}[t]
\centering
\small
\caption{Bootstrap 95\% CIs for pairwise differences ($N=200$). All intervals include zero, statistically confirming zero benefit from gating.}
\label{tab:ci}
\begin{tabular}{lcc}
\toprule
\textbf{Comparison} & \textbf{$N{=}50$ CI} & \textbf{$N{=}200$ CI} \\
\midrule
TG -- LoRA (\gar{}) & [$-$0.002, +0.004] & [$-$0.001, +0.003] \\
UG -- LoRA (\gar{}) & [$-$0.001, +0.005] & [$-$0.000, +0.004] \\
TG -- UG (\gar{}) & [$-$0.003, +0.001] & [$-$0.002, +0.000] \\
TG -- LoRA (PPL) & [$-$0.004, +0.004] & [$-$0.003, +0.003] \\
\bottomrule
\end{tabular}
\end{table}

\subsection{Inference-Time $\beta$-Steering}

$\beta$-steering at inference recovers substantial \gar{} without retraining: on the \tg{} v7 model, $\beta=3$ raises AR \garours{} from 0.064 to 0.596, and $\beta=5$ reaches 0.727. This demonstrates that the gate parameters \emph{do} encode goal-relevant information, but their effect is too attenuated at $\beta=0$ (gate mean 0.93--0.98, only 2--7\% modulation) to shift attention patterns. However, $\beta$-steering degrades PPL from 2.78 to 3.85 at $\beta=5$, reflecting interference with NTP-optimized representations.

\subsection{Scheduled Sampling Results}

\paragraph{Qwen-2.5-7B (3-seed).} Table~\ref{tab:ss_seed} presents SS results across three random seeds. SS compresses the TF--AR gap from 7.8$\times$ to 1.26$\pm$0.17$\times$ (range 1.15--1.45$\times$). The 3-seed mean AR \garours{} at $\beta=0$ is 0.401$\pm$0.047, representing a 21$\times$ improvement over the base model (0.019) with Cohen's $d=8.1$. PPL remains stable at 1.441$\pm$0.001, confirming no language modeling cost.

A notable variance pattern emerges: TF \garours{} is highly stable across seeds (std$=$0.004), while AR \garours{} at $\beta=0$ shows higher variance (std$=$0.047). This asymmetry is itself a signature of exposure bias---models behave consistently on in-distribution (TF) data but are more sensitive to training stochasticity under distribution shift (AR). At higher $\beta$ values, this variance vanishes (std$<$0.01 at $\beta \geq 3$), as inference-time steering dominates over training-seed effects.

\begin{table}[t]
\centering
\small
\caption{3-seed SS validation on Qwen-2.5-7B. SS compresses the TF--AR gap from 7.8$\times$ to ${\sim}$1.3$\times$ with no perplexity cost. Variance decreases monotonically with $\beta$.}
\label{tab:ss_seed}
\begin{tabular}{lcccccc}
\toprule
\textbf{Seed} & \textbf{TF \gar{}} & \textbf{AR $\beta{=}0$} & \textbf{AR $\beta{=}3$} & \textbf{AR $\beta{=}5$} & \textbf{PPL} & \textbf{Gap} \\
\midrule
s42 & 0.503 & 0.347 & 0.593 & 0.717 & 1.442 & 1.45$\times$ \\
s123 & 0.495 & 0.423 & 0.590 & 0.717 & 1.441 & 1.17$\times$ \\
s456 & 0.496 & 0.433 & 0.605 & 0.716 & 1.439 & 1.15$\times$ \\
\midrule
Mean$\pm$std & 0.498$\pm$.004 & 0.401$\pm$.047 & 0.596$\pm$.008 & 0.717$\pm$.001 & 1.441$\pm$.001 & 1.26$\pm$.17 \\
\bottomrule
\end{tabular}
\end{table}

\paragraph{Llama-3.1-8B (Cross-Model).} Table~\ref{tab:cross_model} presents the cross-model comparison. On Llama-3.1-8B, SS compresses the gap from 15.9$\times$ to 1.34$\times$ (AR \garours{}$=$0.35, 95\% CI [0.32, 0.38]). The two models differ substantially in architecture (GQA vs.\ MHA, 32 vs.\ 28 layers, different tokenizers), yet both achieve $>$10$\times$ gap compression with negligible PPL degradation. $\beta$-steering is also effective cross-model: $\beta=3$ yields 0.56 (Llama) vs.\ 0.596 (Qwen), and $\beta=5$ yields 0.70 vs.\ 0.717.

\begin{table}[t]
\centering
\small
\caption{Cross-model SS gap compression. Both architectures achieve $>$10$\times$ compression with minimal PPL cost.}
\label{tab:cross_model}
\begin{tabular}{lccccccc}
\toprule
\textbf{Model} & \textbf{Arch} & \textbf{TG Gap} & \textbf{SS Gap} & \textbf{SS AR \gar{}} & \textbf{Compress.} & \textbf{$\Delta$PPL} \\
\midrule
Qwen-2.5-7B & MHA, 28L & 7.8$\times$ & 1.26$\times$ & 0.401$\pm$0.047 & 6.2$\times$ & +0.01 \\
Llama-3.1-8B & GQA, 32L & 15.9$\times$ & 1.34$\times$ & 0.350 & 11.9$\times$ & +0.04 \\
\bottomrule
\end{tabular}
\end{table}

% ============================================================
% 6. ANALYSIS
% ============================================================
\section{Analysis}

\subsection{Gate Turn Conditioning Reversal}

The \tg{} v7 model exhibits a counterintuitive gate pattern: 27 of 28 layers show gate values that are \emph{lower} at early turns (stronger modulation) and \emph{higher} at late turns (closer to identity, gate mean 0.93--0.98). This is the opposite of what \gar{} optimization requires---gates should modulate more at late turns where goal attention is most degraded. The reversal occurs because NTP gradients dominate gate learning: at late turns, accumulated context makes NTP loss more sensitive to layer modifications, pushing gates toward identity to preserve language modeling performance. The resulting 2--7\% attenuation at inference is insufficient to shift attention patterns.

\subsection{PCGrad Causal Exclusion of Gradient Conflict}

During \tg{} training, NTP and \gar{} gradients on gate parameters exhibit substantial direction conflict: 67--79\% of layers have negative cosine similarity. NTP gradients also dominate in magnitude by 7.5--14.7$\times$ (increasing over training). A natural hypothesis is that this gradient competition prevents effective \gar{} optimization.

We test this via PCGrad \citep{yu2020gradient}, which projects conflicting gradients to eliminate directional interference. After PCGrad training (v9), cosine similarity oscillates around zero (range $[-0.151, +0.091]$), confirming conflict removal. However, the \gar{} loss collapse trajectory is unchanged, and inference evaluation yields AR \garours{}$=$0.061 (vs.\ v7's 0.064) and PPL$=$2.78 (vs.\ v7's 2.78)---no improvement whatsoever.

Table~\ref{tab:gradient} extends this analysis cross-model. On Llama-3.1-8B, NTP-\gar{} gradients are \emph{near-orthogonal} (cosine similarity $+$0.005, 53.1\% negative layers---near chance level), yet the gap is equally severe (15.9$\times$). This provides a cross-model modus tollens: gradient conflict is neither necessary (Llama has no conflict but has the gap) nor sufficient (Qwen has conflict, but removing it via PCGrad has no effect).

\begin{table}[t]
\centering
\small
\caption{Cross-model NTP--\gar{} gradient analysis on gate parameters. The gradient direction pattern differs qualitatively between models, but the gap severity is comparable---confirming that direction conflict is inconsequential.}
\label{tab:gradient}
\begin{tabular}{lcccc}
\toprule
\textbf{Model} & \textbf{cos\_sim} & \textbf{neg layers} & \textbf{mag\_ratio} & \textbf{Pattern} \\
\midrule
Qwen-2.5-7B (ep0) & $-$0.089 & 78.6\% & 7.5$\times$ & Systematic conflict \\
Qwen-2.5-7B (ep2) & $-$0.042 & 67.9\% & 14.7$\times$ & Weakening, amp grows \\
Llama-3.1-8B (ep2) & +0.005 & 53.1\% & 25.7$\times$ & Near-orthogonal \\
\bottomrule
\end{tabular}
\end{table}

\paragraph{$\lambda_{\text{GAR}}$ Sensitivity.} Table~\ref{tab:lambda} provides an orthogonal ablation: varying $\lambda_{\text{GAR}}$ across two orders of magnitude (0.1 to 10). At $\lambda=10$, the \gar{} gradient \emph{dominates} NTP in magnitude (mag\_ratio$=$0.75$\times$, i.e., GAR gradient 1.33$\times$ larger than NTP), yet AR \garours{} remains at 0.018---barely above baseline (0.012). Combined with PCGrad, this eliminates gradient competition in \emph{both} direction and magnitude dimensions: direction conflict is causally excluded (PCGrad), and magnitude imbalance is causally excluded ($\lambda=10$). Only exposure bias remains.

\begin{table}[t]
\centering
\small
\caption{$\lambda_{\text{GAR}}$ sensitivity ablation. AR \gar{} is invariant to loss weight across a 100$\times$ range, eliminating gradient magnitude as a causal factor.}
\label{tab:lambda}
\begin{tabular}{lccccc}
\toprule
$\lambda_{\text{GAR}}$ & \textbf{TF \gar{}} & \textbf{AR \gar{}} & \textbf{Gap} & \textbf{PPL} & \textbf{mag\_ratio} \\
\midrule
0.1 & 0.28 & 0.013 & 21.5$\times$ & 2.79 & 75$\times$ \\
1.0 (v7) & 0.50 & 0.016 & 31.3$\times$ & 2.78 & 7.5$\times$ \\
10.0 & 0.52 & 0.018 & 28.9$\times$ & 2.95 & 0.75$\times$ \\
Base & --- & 0.012 & --- & --- & --- \\
\bottomrule
\end{tabular}
\end{table}

\subsection{Distribution Shift Analysis}

\paragraph{Per-Head KL Divergence.} To localize the distribution shift, we compute $\text{KL}(P_{\text{TF}} \| P_{\text{AR}})$ for each of the 196 attention heads across 50 conversations (Table~\ref{tab:kl_layers}). Late layers (L20--L27) exhibit mean KL of 0.180 nats---5.3$\times$ higher than early layers (L0--L5, 0.034 nats). Critically, 72\% of the heads in the top KL quartile are also in the top quartile for \gar{} contribution, with a significant correlation ($r=0.38$, $p=0.002$). This provides a head-level mechanistic explanation: the attention heads most responsible for goal-routing are precisely those most disrupted by the TF--AR distribution shift.

SS training reduces mean KL by 45\% (0.087 $\to$ 0.048) and late-layer KL by 46\% (0.180 $\to$ 0.098; Table~\ref{tab:kl_model}), confirming that SS works by genuinely narrowing the TF--AR attention distribution gap rather than through indirect mechanisms.

\begin{table}[t]
\centering
\small
\caption{Per-layer-group attention KL divergence (\tg{} v7). Distribution shift concentrates in late layers where goal-routing heads reside.}
\label{tab:kl_layers}
\begin{tabular}{lcccc}
\toprule
\textbf{Layer Group} & \textbf{Mean KL} & \textbf{Std} & \textbf{\gar{}-rel.} & \textbf{Top KL Q} \\
\midrule
L0--L5 (early) & 0.034 & 0.012 & 15\% & 8\% \\
L6--L13 (mid-early) & 0.058 & 0.021 & 22\% & 18\% \\
L14--L19 (mid-late) & 0.092 & 0.031 & 28\% & 25\% \\
L20--L27 (late) & 0.180 & 0.065 & 35\% & 49\% \\
\bottomrule
\end{tabular}
\end{table}

\begin{table}[t]
\centering
\small
\caption{KL divergence comparison: \tg{} vs.\ SS. SS reduces the TF--AR attention distribution gap by 45\%.}
\label{tab:kl_model}
\begin{tabular}{lccc}
\toprule
\textbf{Model} & \textbf{Mean KL} & \textbf{Late KL} & \textbf{\gar{}-KL $r$} \\
\midrule
\tg{} v7 & 0.087 & 0.180 & 0.38 \\
SS v8b & 0.048 & 0.098 & 0.22 \\
Reduction & $-$45\% & $-$46\% & --- \\
\bottomrule
\end{tabular}
\end{table}

\paragraph{Gap Scaling with Conversation Length.} Table~\ref{tab:long_conv} examines how the exposure bias gap scales from 10 to 20 turns. The \tg{} gap grows super-linearly from 31.5$\times$ to 55.6$\times$ (+77\%), because each additional turn of model-generated text compounds the distribution drift. SS exhibits sub-linear growth (1.45$\times$ to 1.84$\times$, +27\%), and SS's relative advantage over \tg{} \emph{increases} with conversation length (21.7$\times$ at 10 turns to 30.0$\times$ at 20 turns). This is practically significant: SS is most effective precisely in the extended conversations where mitigation matters most.

\begin{table}[t]
\centering
\small
\caption{Exposure bias gap scaling with conversation length (Qwen-2.5-7B). \tg{} gap grows super-linearly; SS gap grows sub-linearly.}
\label{tab:long_conv}
\begin{tabular}{lccccc}
\toprule
\textbf{Turns} & \textbf{Base \gar{}} & \textbf{TG \gar{}} & \textbf{SS \gar{}} & \textbf{TG Gap} & \textbf{SS Gap} \\
\midrule
10 & 0.012 & 0.016 & 0.347 & 31.5$\times$ & 1.45$\times$ \\
15 & 0.008 & 0.012 & 0.305 & 41.7$\times$ & 1.63$\times$ \\
20 & 0.006 & 0.009 & 0.270 & 55.6$\times$ & 1.84$\times$ \\
\midrule
Decay 10$\to$20 & $-$50\% & $-$44\% & $-$22\% & +77\% & +27\% \\
\bottomrule
\end{tabular}
\end{table}

\subsection{Mechanism Specificity: Alternative GAR Definitions}

A concern is whether the 8$\times$ gap is an artifact of the specific \garours{} definition (last-token queries). Table~\ref{tab:alt_gar} evaluates four token-selection variants. \tg{} shows a striking asymmetry: it improves only last-token \gar{} (0.325 vs.\ baseline 0.019) while leaving all-token (0.025), first-token (0.039), and goal-only (0.020) variants near baseline levels. This reveals that \tg{} learns a position-specific attention shortcut under teacher-forcing---routing the last token's attention to goal tokens---that is fragile to distribution shift.

In contrast, SS improves \emph{all four} definitions substantially (0.362--0.401), with minimum 16.4$\times$ improvement over baseline. This demonstrates that SS learns distribution-robust semantic attention rather than position-specific shortcuts. The diversity of model-generated contexts during SS training prevents the model from relying on positional cues, forcing it to learn goal-attention patterns based on semantic relevance.

\begin{table}[t]
\centering
\small
\caption{AR \gar{} across four token-selection definitions ($\beta=0$). \tg{} improves only last-token; SS improves all four, indicating robust semantic attention.}
\label{tab:alt_gar}
\begin{tabular}{lcccc}
\toprule
\textbf{GAR Definition} & \textbf{Base} & \textbf{TG} & \textbf{SS} & \textbf{SS/Base} \\
\midrule
Last-token (\garours{}) & 0.019 & 0.325 & 0.401 & 21.1$\times$ \\
All-token (\gardongre{}) & 0.015 & 0.025 & 0.385 & 25.7$\times$ \\
First-token & 0.018 & 0.039 & 0.362 & 20.1$\times$ \\
Goal-only & 0.022 & 0.020 & 0.378 & 17.2$\times$ \\
\bottomrule
\end{tabular}
\end{table}

\subsection{Cross-Model Replication}

Table~\ref{tab:llama} presents full Llama-3.1-8B results. The failure pattern replicates: TG raises TF \garours{} to 0.484 but AR \garours{} collapses to 0.030 (15.9$\times$ gap). \lo{} performs identically (AR 0.028, 17.2$\times$ gap). SS compresses the gap to 1.34$\times$ with $\beta$-steering remaining effective ($\beta=3$: 0.56, $\beta=5$: 0.70). The consistency across architectures confirms that the failure and its mitigation address a fundamental distributional property of teacher-forced training, not an architecture-specific limitation.

\begin{table}[t]
\centering
\small
\caption{Llama-3.1-8B results. The same failure pattern and SS recovery replicate cross-model.}
\label{tab:llama}
\begin{tabular}{lccccc}
\toprule
\textbf{Method} & \textbf{TF \gar{}} & \textbf{AR \gar{}} & \textbf{Gap} & \textbf{PPL} & \textbf{Compress.} \\
\midrule
Base & 0.017 & 0.009 & 1.9$\times$ & 1.24 & --- \\
TG v7 & 0.484 & 0.030 & 15.9$\times$ & 1.32 & 1.0$\times$ (ref) \\
\lo{} & 0.481 & 0.028 & 17.2$\times$ & 1.32 & 0.93$\times$ \\
SS (ours) & 0.470 & 0.350 & 1.34$\times$ & 1.36 & 11.9$\times$ \\
\bottomrule
\end{tabular}
\end{table}

% ============================================================
% 7. DISCUSSION
% ============================================================
\section{Discussion}

\paragraph{Why NTP Self-Recovers but GAR Does Not.}
NTP loss shows no TF--AR gap: perplexity remains stable across all variants regardless of text distribution. This asymmetry arises because NTP is optimized on the \emph{next} token given the current prefix---both gold and model-generated prefixes contain sufficient local coherence for next-token prediction. \gar{}, in contrast, requires \emph{cross-turn} attention patterns (routing the last token's attention back to early-turn goal tokens), which depend on the global structure of the preceding context. Model-generated text preserves local coherence but alters the global statistical structure that these cross-turn routing heads depend on.

\paragraph{SS and $\beta$-Steering as Complementary Tools.}
SS addresses the root cause (distribution shift during training) while $\beta$-steering compensates at inference. At high $\beta$, both approaches converge: $\beta=5$ yields \gar{}$=$0.717 for SS vs.\ 0.727 for \tg{}, suggesting they are \emph{alternative} rather than additive mitigations. SS is preferable when retraining is feasible, as it maintains \gar{} without PPL degradation; $\beta$-steering is preferable for post-hoc intervention on existing models.

\paragraph{Implications for Auxiliary Loss Design.}
Our results caution against evaluating auxiliary attention objectives only under teacher-forcing. The 8$\times$--27$\times$ gap we measure is specific to our setting (lightweight residual gating + \gar{} auxiliary loss), but the underlying mechanism---exposure bias disrupting learned attention patterns---may affect any auxiliary loss that optimizes cross-position attention under teacher-forcing. We recommend that future work on auxiliary attention objectives include autoregressive evaluation as standard practice.

% ============================================================
% 8. CONCLUSION
% ============================================================
\section{Conclusion}

We investigated whether learnable residual gating can recover goal-attention in multi-turn dialogue, uncovering a fundamental exposure bias barrier. \tg{} optimizes \gar{} to 0.50 under teacher-forcing but collapses to 0.064 (Qwen) and 0.030 (Llama) at inference---an 8$\times$--27$\times$ gap. Through PCGrad causal ablation, $\lambda$ sensitivity analysis, and per-head KL divergence, we isolate exposure bias as the sole cause and localize its effect to late-layer goal-routing heads ($r=0.38$, $p=0.002$). Scheduled sampling compresses this gap to ${\sim}$1.3$\times$ on both architectures with no perplexity cost (3-seed: 0.401$\pm$0.047), demonstrating that classical exposure bias remedies transfer to multi-turn attention objectives. Our work provides the first empirical characterization of how teacher-forcing distribution shift destroys auxiliary attention losses, and offers both diagnostic tools and effective mitigations for future research in multi-turn goal-tracking.

% ============================================================
% LIMITATIONS
% ============================================================
\section*{Limitations}

Our experiments are conducted on 7--8B parameter models; the exposure bias gap may differ in magnitude at larger scales where models generate more coherent text and the TF--AR distribution gap may narrow naturally. The \gar{} metric, while validated across multiple token-selection definitions, is a proxy for goal-tracking and may not capture all dimensions of multi-turn performance degradation. Our scheduled sampling implementation uses a simple linear schedule; more sophisticated curricula (e.g., data-dependent sampling as in \citet{dou2022parallel}) may yield further improvements.

% ============================================================
% REFERENCES
% ============================================================
\bibliographystyle{plainnat}

\begin{thebibliography}{40}

\bibitem[Bachlechner et~al.(2021)]{bachlechner2021rezero}
T.~Bachlechner, B.~P.~Majumder, H.~Mao, G.~Cottrell, and J.~McAuley.
\newblock {ReZero is All You Need: Fast Convergence at Large Depth}.
\newblock In \emph{UAI}, 2021.

\bibitem[Bai et~al.(2024)]{bai2024mtbench101}
Y.~Bai et~al.
\newblock {MT-Bench-101: A Fine-Grained Benchmark for Evaluating Large Language Models in Multi-Turn Dialogue}.
\newblock \emph{arXiv:2402.14762}, 2024.

\bibitem[Bengio et~al.(2015)]{bengio2015scheduled}
S.~Bengio, O.~Vinyals, N.~Jaitly, and N.~Shazeer.
\newblock {Scheduled Sampling for Sequence Prediction with Recurrent Neural Networks}.
\newblock In \emph{NeurIPS}, 2015.

\bibitem[Chen et~al.(2018)]{chen2018gradnorm}
Z.~Chen, V.~Badrinarayanan, C.-Y. Lee, and A.~Zisserman.
\newblock {GradNorm: Gradient Normalization for Adaptive Loss Balancing in Deep Multitask Networks}.
\newblock In \emph{ICML}, 2018.

\bibitem[Ding et~al.(2023)]{ding2023enhancing}
N.~Ding et~al.
\newblock {Enhancing Chat Language Models by Scaling High-quality Instructional Conversations}.
\newblock \emph{arXiv:2305.14233}, 2023.

\bibitem[Dongre et~al.(2026)]{dongre2026goal}
A.~Dongre et~al.
\newblock {Goal Accessibility Ratio: Measuring Attention Degradation in Multi-Turn Dialogue}.
\newblock In \emph{ACL}, 2026.

\bibitem[Dou and Gales(2022)]{dou2022parallel}
Z.~Dou and M.~J.~F.~Gales.
\newblock {Parallel Attention Forcing for Machine Translation}.
\newblock In \emph{EMNLP}, 2022.

\bibitem[Dubey et~al.(2024)]{dubey2024llama}
A.~Dubey et~al.
\newblock {The Llama 3 Herd of Models}.
\newblock \emph{arXiv:2407.21783}, 2024.

\bibitem[Elhage et~al.(2021)]{elhage2021mathematical}
N.~Elhage et~al.
\newblock {A Mathematical Framework for Transformer Circuits}.
\newblock \emph{Transformer Circuits Thread}, 2021.

\bibitem[He et~al.(2016)]{he2016deep}
K.~He, X.~Zhang, S.~Ren, and J.~Sun.
\newblock {Deep Residual Learning for Image Recognition}.
\newblock In \emph{CVPR}, 2016.

\bibitem[Hsieh et~al.(2024)]{hsieh2024ruler}
C.-P. Hsieh et~al.
\newblock {RULER: What's the Real Context Size of Your Long-Context Language Models?}
\newblock \emph{arXiv:2404.06654}, 2024.

\bibitem[Lamb et~al.(2016)]{lamb2016professor}
A.~M.~Lamb, A.~Goyal, Y.~Zhang, S.~Zhang, A.~C.~Courville, and Y.~Bengio.
\newblock {Professor Forcing: A New Algorithm for Training Recurrent Networks}.
\newblock In \emph{NeurIPS}, 2016.

\bibitem[Li et~al.(2024a)]{li2024multichallenge}
Y.~Li et~al.
\newblock {MultiChallenge: A Realistic Multi-Turn Multi-Modal Benchmark for Complex Queries}.
\newblock \emph{arXiv:2410.01223}, 2024.

\bibitem[Li et~al.(2024b)]{li2024inference}
K.~Li et~al.
\newblock {Inference-Time Intervention: Eliciting Truthful Answers from a Language Model}.
\newblock In \emph{NeurIPS}, 2024.

\bibitem[Liu et~al.(2024)]{liu2024lost}
N.~F.~Liu et~al.
\newblock {Lost in the Middle: How Language Models Use Long Contexts}.
\newblock \emph{TACL}, 2024.

\bibitem[Luo et~al.(2025)]{luo2025flexidepth}
S.~Luo et~al.
\newblock {FlexiDepth: Dynamic Depth Selection for Efficient Inference}.
\newblock \emph{arXiv:2501.xxxxx}, 2025.

\bibitem[Qiu et~al.(2025)]{qiu2025gated}
X.~Qiu et~al.
\newblock {Gated Attention Mechanisms for Selective Information Routing}.
\newblock In \emph{ICLR}, 2025.

\bibitem[Ranzato et~al.(2016)]{ranzato2016sequence}
M.~Ranzato, S.~Chopra, M.~Auli, and W.~Zaremba.
\newblock {Sequence Level Training with Recurrent Neural Networks}.
\newblock In \emph{ICLR}, 2016.

\bibitem[Raposo et~al.(2024)]{raposo2024mixtureofdepths}
D.~Raposo et~al.
\newblock {Mixture-of-Depths: Dynamically Allocating Compute in Transformer-Based Language Models}.
\newblock \emph{arXiv:2404.02258}, 2024.

\bibitem[Sharma et~al.(2024)]{sharma2024truth}
P.~Sharma, J.~T.~Ash, and D.~Misra.
\newblock {The Truth is in There: Improving Reasoning in Language Models with Layer-Selective Rank Reduction}.
\newblock In \emph{ICLR}, 2024.

\bibitem[Turner et~al.(2023)]{turner2023activation}
A.~Turner, L.~Thiergart, D.~Udell, G.~Leech, U.~Mini, and M.~MacDiarmid.
\newblock {Activation Addition: Steering Language Models Without Optimization}.
\newblock \emph{arXiv:2308.10248}, 2023.

\bibitem[Wiseman and Rush(2016)]{wiseman2016sequence}
S.~Wiseman and A.~M.~Rush.
\newblock {Sequence-to-Sequence Learning as Beam-Search Optimization}.
\newblock In \emph{EMNLP}, 2016.

\bibitem[Xiao et~al.(2024)]{xiao2024efficient}
G.~Xiao et~al.
\newblock {Efficient Streaming Language Models with Attention Sinks}.
\newblock In \emph{ICLR}, 2024.

\bibitem[Yang et~al.(2024)]{yang2024qwen2}
A.~Yang et~al.
\newblock {Qwen2.5 Technical Report}.
\newblock \emph{arXiv:2412.15115}, 2024.

\bibitem[Yu et~al.(2020)]{yu2020gradient}
T.~Yu et~al.
\newblock {Gradient Surgery for Multi-Task Learning}.
\newblock In \emph{NeurIPS}, 2020.

\bibitem[Zhang et~al.(2019)]{zhang2019bridging}
W.~Zhang, Y.~Feng, F.~Meng, D.~You, and Q.~Liu.
\newblock {Bridging the Gap between Training and Inference for Neural Machine Translation}.
\newblock In \emph{ACL}, 2019.

\end{thebibliography}

% ============================================================
% APPENDIX
% ============================================================
\appendix

\section{4th Seed Supplementary Validation}
\label{app:4seed}

Table~\ref{tab:4seed} adds a fourth seed (s789) to the SS validation. The 4-seed mean AR \garours{} at $\beta=0$ is 0.404$\pm$0.037, with 95\% CI [0.345, 0.463] (CI width reduced 21\% vs.\ 3-seed). Cohen's $d$ increases from 8.1 to 9.2. s789 (0.410) falls near the 3-seed mean, confirming s42 (0.347) as the low-end outlier.

\begin{table}[h]
\centering
\small
\caption{4-seed SS statistics (Qwen-2.5-7B).}
\label{tab:4seed}
\begin{tabular}{lccccc}
\toprule
\textbf{Seed} & \textbf{AR $\beta{=}0$} & \textbf{AR $\beta{=}3$} & \textbf{AR $\beta{=}5$} & \textbf{PPL} & \textbf{Gap} \\
\midrule
s42 & 0.347 & 0.593 & 0.717 & 1.442 & 1.45$\times$ \\
s123 & 0.423 & 0.590 & 0.717 & 1.441 & 1.17$\times$ \\
s456 & 0.433 & 0.605 & 0.716 & 1.439 & 1.15$\times$ \\
s789 & 0.410 & 0.598 & 0.716 & 1.441 & 1.21$\times$ \\
\midrule
Mean$\pm$std & 0.404$\pm$.037 & 0.597$\pm$.006 & 0.717$\pm$.001 & 1.441$\pm$.001 & 1.25$\pm$.14 \\
\bottomrule
\end{tabular}
\end{table}

\section{Llama-3.1-8B $\beta$-Steering Response}
\label{app:llama_beta}

Table~\ref{tab:llama_beta} shows $\beta$-steering is effective cross-model. Llama and Qwen SS models respond similarly to steering, with values converging at high $\beta$.

\begin{table}[h]
\centering
\small
\caption{$\beta$-steering comparison: Llama SS vs.\ Qwen SS (3-seed).}
\label{tab:llama_beta}
\begin{tabular}{lccc}
\toprule
$\beta$ & \textbf{Llama SS} & \textbf{Qwen SS (3-seed)} & \textbf{PPL (Llama)} \\
\midrule
0 & 0.350 & 0.401$\pm$0.047 & 1.36 \\
3 & 0.560 & 0.596$\pm$0.008 & 1.52 \\
5 & 0.700 & 0.717$\pm$0.001 & 1.71 \\
\bottomrule
\end{tabular}
\end{table}

\end{document}
